% GENERATED FILE. Edit canonical lessons, then run npm run generate:books.
% canonical-source-hash: fnv1a64:540b5ee72c77b0b5
% canonical-lessons: FR-C510-patin, FR-C510-patiner, FR-C510-luge, FR-C510-patinoire, FR-C510-piste

\chapter{Glace et neige}
\label{ch:fr-glace-et-neige}
\hlchaptermodality{510}

\begin{chapteropening}
\textbf{By the end of this chapter:} \emph{I can talk about skating and sledging.}

Complete the last lesson of chapter 510: I can talk about skating and sledging.
\end{chapteropening}

\section[le patin]{le patin — skate}

\label{lesson:FR-C510-patin}



\begin{quote}
Before the new one: say the French for breaststroke, then the French for to row.
\end{quote}

\subsection*{You'll want to know: le patin}
\textbf{le patin} — \textquotedblleft{}skate\textquotedblright{}.

\begin{quote}
J'ai reçu des patins pour mon anniversaire.
\end{quote}

\begin{quote}
\emph{I got skates for my birthday.}
\end{quote}

\emph{Le patin à glace} — ice skating; \emph{le patin à roulettes} — roller skating.

\begin{cousinweb}
\emph{le patin} is built on \emph{patte} \textquotedblleft{}paw, foot\textquotedblright{} with \emph{-in} (EtymDB); the origin of \emph{patte} is not settled.
\end{cousinweb}

\begin{grammarlens}[title={What you've built}]
One more word for telling, explaining and describing at B1.
\end{grammarlens}

\subsection*{Your turn}
\begin{itemize}
\raggedright
  \item \emph{Say it:} \emph{le patin}
  \item \emph{Say it:} \emph{le patin}, once more
  \item \emph{Say it:} a sentence of your own with \emph{le patin}
  \item \emph{Recall:} say the French for breaststroke, then the French for to row, then say \emph{le patin} again
\end{itemize}

\begin{tcolorbox}[breakable,colback=teal!4,colframe=teal!35!black,title={Before you move on}]
What does le patin mean? (\textquotedblleft{}Skate\textquotedblright{}.) Say it once more.
\end{tcolorbox}
\section[patiner]{patiner — to skate}

\label{lesson:FR-C510-patiner}



\begin{quote}
Before the new one: say the French for to row, then the French for skate.
\end{quote}

\subsection*{You'll want to know: patiner}
\textbf{patiner} — \textquotedblleft{}to skate\textquotedblright{}.

\begin{quote}
L'hiver, nous patinons sur le lac gelé.
\end{quote}

\begin{quote}
\emph{In winter we skate on the frozen lake.}
\end{quote}

Of wheels: to spin: \emph{la voiture patine}.

\begin{cousinweb}
\emph{patiner} is built on \emph{patin} \textquotedblleft{}skate\textquotedblright{}, from \emph{patte} \textquotedblleft{}paw\textquotedblright{}, whose origin is not settled.
\end{cousinweb}

\begin{grammarlens}[title={What you've built}]
One more word for telling, explaining and describing at B1.
\end{grammarlens}

\subsection*{Your turn}
\begin{itemize}
\raggedright
  \item \emph{Say it:} \emph{patiner}
  \item \emph{Say it:} \emph{patiner}, once more
  \item \emph{Say it:} a sentence of your own with \emph{patiner}
  \item \emph{Recall:} say the French for to row, then the French for skate, then say \emph{patiner} again
\end{itemize}

\begin{tcolorbox}[breakable,colback=teal!4,colframe=teal!35!black,title={Before you move on}]
What does patiner mean? (\textquotedblleft{}To skate\textquotedblright{}.) Say it once more.
\end{tcolorbox}
\section[la luge]{la luge — sledge, sled}

\label{lesson:FR-C510-luge}



\begin{quote}
Before the new one: say the French for skate, then the French for to skate.
\end{quote}

\subsection*{You'll want to know: la luge}
\textbf{la luge} — \textquotedblleft{}sledge, sled\textquotedblright{}.

\begin{quote}
Les enfants font de la luge sur la colline.
\end{quote}

\begin{quote}
\emph{The children are sledging on the hill.}
\end{quote}

\emph{Faire de la luge} — to go sledging. In Québec \emph{la traîne sauvage}.

\begin{cousinweb}
\emph{la luge} came into French from the dialects of the Alps. It is usually traced to a Gaulish word; the origin is not fully settled.
\end{cousinweb}

\begin{grammarlens}[title={What you've built}]
One more word for telling, explaining and describing at B1.
\end{grammarlens}

\subsection*{Your turn}
\begin{itemize}
\raggedright
  \item \emph{Say it:} \emph{la luge}
  \item \emph{Say it:} \emph{la luge}, once more
  \item \emph{Say it:} a sentence of your own with \emph{la luge}
  \item \emph{Recall:} say the French for skate, then the French for to skate, then say \emph{la luge} again
\end{itemize}

\begin{tcolorbox}[breakable,colback=teal!4,colframe=teal!35!black,title={Before you move on}]
What does la luge mean? (\textquotedblleft{}Sledge, sled\textquotedblright{}.) Say it once more.
\end{tcolorbox}
\section[la patinoire]{la patinoire — ice rink}

\label{lesson:FR-C510-patinoire}



\begin{quote}
Before the new one: say the French for to skate, then the French for sledge.
\end{quote}

\subsection*{You'll want to know: la patinoire}
\textbf{la patinoire} — \textquotedblleft{}ice rink\textquotedblright{}.

\begin{quote}
Le samedi, la patinoire est pleine de monde.
\end{quote}

\begin{quote}
\emph{On Saturdays the ice rink is crowded.}
\end{quote}

Informally, a very icy road: \emph{la route est une patinoire}.

\begin{cousinweb}
\emph{la patinoire} is built on \emph{patiner} \textquotedblleft{}to skate\textquotedblright{} with \emph{-oire}, the ending for a place.
\end{cousinweb}

\begin{grammarlens}[title={What you've built}]
One more word for telling, explaining and describing at B1.
\end{grammarlens}

\subsection*{Your turn}
\begin{itemize}
\raggedright
  \item \emph{Say it:} \emph{la patinoire}
  \item \emph{Say it:} \emph{la patinoire}, once more
  \item \emph{Say it:} a sentence of your own with \emph{la patinoire}
  \item \emph{Recall:} say the French for to skate, then the French for sledge, then say \emph{la patinoire} again
\end{itemize}

\begin{tcolorbox}[breakable,colback=teal!4,colframe=teal!35!black,title={Before you move on}]
What does la patinoire mean? (\textquotedblleft{}Ice rink\textquotedblright{}.) Say it once more.
\end{tcolorbox}
\section[la piste]{la piste — slope, run; track}

\label{lesson:FR-C510-piste}



\begin{quote}
Before the new one: say the French for sledge, then the French for ice rink.
\end{quote}

\subsection*{You'll want to know: la piste}
\textbf{la piste} — \textquotedblleft{}slope, run; track\textquotedblright{}.

\begin{quote}
Cette piste noire est trop difficile pour moi.
\end{quote}

\begin{quote}
\emph{This black run is too difficult for me.}
\end{quote}

Also a running track, a dance floor, or a clue: \emph{une bonne piste}.

\begin{cousinweb}
\emph{la piste} was borrowed from Italian \emph{pista} \textquotedblleft{}a beaten track\textquotedblright{} (EtymDB), built on \emph{pistare} \textquotedblleft{}to pound\textquotedblright{}, which goes back to Latin \emph{pinsere} \textquotedblleft{}to pound\textquotedblright{} (Lewis \& Short).
\end{cousinweb}

\begin{grammarlens}[title={What you've built}]
One more word for telling, explaining and describing at B1.
\end{grammarlens}

\subsection*{Your turn}
\begin{itemize}
\raggedright
  \item \emph{Say it:} \emph{la piste}
  \item \emph{Say it:} \emph{la piste}, once more
  \item \emph{Say it:} a sentence of your own with \emph{la piste}
  \item \emph{Recall:} say the French for sledge, then the French for ice rink, then say \emph{la piste} again
\end{itemize}

\begin{tcolorbox}[breakable,colback=teal!4,colframe=teal!35!black,title={Before you move on}]
What does la piste mean? (\textquotedblleft{}Slope, run; track\textquotedblright{}.) Say it once more.
\end{tcolorbox}
